% NOETHER PAPER 4 — KOREAN TRANSLATION-PRODUCER DRAFT — UNCHECKED
% Unit: T01-U01; current-authority whole lines 3611--3640
% Source slice: 10,273 LF UTF-8 bytes; SHA-256 9A1B0B21FD34C151E7A8B8605B67CE37DEB1F47E8B00C26282F879335FF88591
% Pointer: NOETH-DE-AUTH-v003-20260804; SHA-256 932FEDC1735A41A9CF71D15A6C662A468A4CAD016AE8B3DECDF9A71E8BA7F197
% This is producer translation only: no source/Korean/formula review, compilation, or rendering.

\begin{center}
{\Large\bfseries $n$개 변수의 형식에 관한 불변식론\footnote{잘츠부르크 자연과학자 대회에서 발표한 짧은 보고를 참조하라. \emph{Jahresbericht d. deutsch. Math.-Vereinigung}, Bd.\ 19, S.\ 101.}}\par
\vspace{0.5em}
에를랑겐의 \emph{Emmy Noether} 양.\par
\medskip
\rule{0.22\linewidth}{0.4pt}\par
\vspace{0.7em}
Journal f.\ d.\ reine u.\ angew.\ Math.\ 139 (1911), S.\ 118--154
\end{center}

\begin{center}
\textbf{서론.}
\end{center}
$n$개 변수의 형식에 관한 사영 불변식론에서는, 이항 및 삼항 영역을 넘어서면 주된 문제, 즉 형식계의 유한성 증명에 뒤따르는 문제들이 거의 다루어지지 않았다. 그것은 형식들 사이의 상호 의존성과 이 연관관계를 다루는 방법에 관한 문제들이다. 이 방법들은 본질적으로 --- 이항 및 삼항 영역에서는 이것이 완전히 증명되어 있듯이 --- Study가 \glqq{}기호적 방법의 기본정리\grqq{}라고 부른 두 정리, 곧 형식의 기호적 표현 가능성에 관한 정리와 기호 항등식에 관한 정리에 기초한다. 뒤의 정리는 정수 불변 구성식들 사이에 성립하는 모든 관계를 소수의 기호 항등식을 차례로 적용하여 얻는다는 것, 따라서 기호적 방법만이 불변식의 영역을 벗어나지 않고 형식들의 연관관계를 파악하게 한다는 것을 말한다.\footnote{E. Study, \emph{Methoden zur Theorie der ternären Formen} (Leipzig, Teubner, 1889). II. \S\ 6. --- 특수 변환군의 불변식론에서도 이 기본정리들이 나타나는 것에 관해서는 Study, \emph{Das Apollonische Problem} (Math. Ann. Bd. 49 [1897]) 및 \emph{Zur Differentialgeometrie der analytischen Kurven} (Transactions of the American Math. Society, Bd. 1 [1909])을 참조하라.}

$n$개 변수로 넘어갈 때 생기는 어려움의 근거는 Gordan이 이미 자신의 에를랑겐 강령에서 지적하였다.\footnote{P. Gordan, \emph{Über das Formensystem binärer Formen}. S. 46. (Leipzig, Teubner, 1875).} 그것은 기호적 방법의 첫째 정리가 그 일반적 형태로는 더 이상 성립하지 않는 데 있다. 잘 알려져 있듯이 $n$항 영역에서는 더 높은 단계의 변수열 $p_\rho$\footnote{표기는 \S\ 1을 참조하라.}가 나타나며, 이와 마찬가지로 더 높은 단계의 기호열 $S^\rho$도 나타난다. 실제로 처음 Clebsch가 했듯이\footnote{A. Clebsch, Über eine Fundamentalaufgabe der Invariantentheorie. (Abh. der Ges. d. Wiss. zu Göttingen, Bd. XVII, 1872).} 임의의 개수의 열 $p_\rho$를 지닌 형식에서 출발하든, F. Deruyts와 A. Capelli의 방식에 따라 서로 동변인 임의의 개수의 열 $x$만을 지닌 형식에서 출발하든, 이 열들에 도달한다.\footnote{A. Capelli, \emph{Lezioni sulla teoria delle forme algebriche} (Napoli 1902), \S\ 22--24 및 거기에 제시된 문헌을 참조하라.} 그런데 열 $p_\rho,S^\rho$를 명시적으로 포함하는 기호적 표현은 존재하지 않는다. 따라서 $p_\rho,S^\rho$를 각각의 열 $x$와 이에 반변인 열 $s$로 분해하고, 이들을 여러 행렬식에 분배해야 한다. 물론 미분 조건(공식 (19.) 참조)을 이용하면 주어진 행렬식 결합식이 오직 열 $p_\rho,S^\rho$에만 의존한다는 성질을 검증할 수 있다. 그러나 이 방법으로는 불변 구성식 전체와 그것을 생성하는 과정 전체를 조망할 수 없다.\footnote{Study에서 유래한, 계산상 매우 유용한 기호법의 한 발전(전달: F. Engel, Ber. d. K. Sächs. Ges. d. Wiss., math.-phys. Kl., 1900, S. 126; 4항 영역에 대해서는 E. Study, \emph{Geometrie der Dynamen}, S. 135)도 여기서 말하는 의미의 명시적 표현은 아니다. 예컨대 그것으로부터 형식을 생성하는 비기호적 미분 과정을 얻기는 어려울 것이다.}

이제 우리는 \S\S\ 2와 6에서 \glqq{}대응 행렬에 관한 정리\grqq{}\footnote{1862년 Graßmann의 확장론 Nr. 112를 제외하면, A. Brill, Über zwei Berührungsprobleme, \S\ 2 (Math. Ann. Bd. 4. 1871)에 처음 나타난다.}를 적용하여, 행렬식에 의한 표현을 \glqq{}행렬곱\grqq{}에 의한 표현으로 바꿈으로써 첫째 기본정리를 일반적인 형태로 다시 얻는 방법을 보인다. \S\ 2에서는 열 $S^\rho,p_\rho$를 명시적으로 포함하는 모든 행렬곱이 기본형의 불변 구성식임을 보인다. \S\ 6에 주어진 역, 즉 모든 불변 구성식이 행렬곱으로 명시적으로 표현된다는 사실은 둘째 기본정리를 전이한 뒤에야(\S\S\ 3--5) 얻어진다.

이 정리는 변수열 $x$만을 포함하는 형식에 대해서는 알려져 있다.\footnote{E. Pascal, Giorn. di mat. 26 (1888), S. 33, 102; Rendic. d. Accad. d. Linc. (4) 4 (1888), S. 119; ib. Mem. (4) 5 (1888), S. 375.} 모든 열 $S^\rho,p_\rho$를 분해 불가능한 것으로 보았을 때의 분해 불가능한 항등식 전체를 세우는 일반 문제는 Study가 정식화하였다.\footnote{\glqq{}Methoden \dots\grqq{}, S. 204, Anmerkung 17).} 그는 동시에 나타나는 항등식의 개수를 $n$항 영역에서 방법의 어려움을 재는 척도로 본다. 이제 항등식 전체를 하나의 공식 (36.)에 모을 수 있으므로, 이 어려움은 적어도 최소한으로 줄어든다. 다만 실제 적용에서는 (36.)의 몇몇 두드러진 특수한 경우가 자주 나타난다. 예컨대 새 열을 대입하지 않고도 명시적 표현을 허용한다는 특징을 지닌 (20.), (21.), 또는 열 $p_\rho$가 겉보기에는 열 $y$들로 분해되어 나타나므로 \glqq{}분해 항등식\grqq{}이라 부르는 공식 (31.)이 그러하다.

이 두 기본정리 위에서 나머지 이론은 쉽게 세울 수 있다. 첫째 정리는 형식을 생성하는 과정, 즉 기호적 수축 과정과 비기호적 미분 과정을 직접 준다.\footnote{\glqq{}Invariante Gebilde von Nullsystemen\grqq{} (Sitzungsber. d. Ak. d. Wiss. Wien, Bd. 97, Abt. IIa, 1888)에서 Mertens는 다른 방식, 곧 직접 일반화할 수 없는 방식으로 4항 수축 과정에 도달하였다. 거기서 나타나는 여섯 열 $S^2$의 교대 불변식은 여기서 대입 (11.)을 한 번 적용하여 생기는 불변식과 짝수 불변식의 한 결합식만큼 다르다. 특수한 한 경우의 4항 수축 과정은 P. Gordan, \emph{Das simultane System von zwei quadratischen quaternären Formen} (Math. Ann. Bd. 56. 1903)에도 나타난다.} 한편 분해 항등식은 이 과정들 가운데 서로 동치인 것을 모두 제거한다(\S\ 6). 더 나아가 미분 과정을 알면 \glqq{}정규형\grqq{}이라는 개념을 $n$항 영역으로 옮길 수 있고, Mertens가 4항 영역에서 제시한 증명법\footnote{F. Mertens, Über invariante Gebilde quaternärer Formen. (Sitzber. d. Ak. d. Wiss. Wien. Bd. 98. Abt. IIa. 1889.)}을 써서 불변 구성식의 한 계를 그와 동치인 정규형의 계로 바꿀 수 있다는 정리를 얻는다(\S\ 7). 나는 이러한 마지막 가정 아래에서, 삼항 수축 과정을 두 기본 수축을 차례로 적용하여 생성할 수 있음을 학위논문에서 보였다.\footnote{Über die Bildung des Formensystems der ternären biquadratischen Form. (Dieses Journal, Bd. 134. 1908.)} 이 정리와 급수 전개 이론에 기초하여 \glqq{}형식열\grqq{} 개념을 도입하고 형식계의 주요 환원정리를 전개하였다. 완전히 마찬가지로 $n$항 영역에서는 수축 과정을 $n-1$개의 기본 수축으로부터 생성할 수 있고(\S\ 8), 환원정리를 세울 수 있다(\S\ 9).

\paragraph{후기.} 잘츠부르크 보고 때 빈의 E. Müller 교수는 이 논문의 바탕을 이루는 행렬곱 계산이 원리상 Graßmann의 확장론 계산법과 동일하다는 점을 내게 알려 주었다.\footnote{Hermann Graßmanns gesammelte mathematische und physikalische Werke. Ersten Bandes zweiter Teil: Die Ausdehnungslehre von 1862. In Gemeinschaft mit H. Graßmann d. Jüngeren herausgegeben v. Fr. Engel (Leipzig, Teubner 1896).} 실제로 Graßmann의 \glqq{}조합곱\grqq{} $[S^{\rho_1}S^{\rho_2}\cdots S^{\rho_i}]$은 이 열들로 주어진 행렬로 볼 수 있다(\S\ 2 참조). 이때 \glqq{}전진곱\grqq{}은 열 자체의 행렬에, \glqq{}후퇴곱\grqq{}은 대응하는 열 $S_{n-\rho_1},S_{n-\rho_2},\ldots,S_{n-\rho_i}$의 행렬에 해당한다. \glqq{}혼합곱\grqq{}에 해당하는 행렬은 여러 열 $S$를 대입 (9.)로 새 열 $P,Q$에 모아 만들게 된다. 여기서 우리의 \glqq{}대응 열\grqq{}은 Graßmann의 \glqq{}보충\grqq{}에, 우리의 \glqq{}행렬곱\grqq{}은 \glqq{}단계수가 0인 혼합곱\grqq{}에 해당한다.

이 대응 아래에서 1862년 확장론 Nr. 113에 주어진 전개는 우리의 공식 (32.)와 쌍대인 공식과 동일해진다. 특수한 경우 $\rho=1$에서 (32.)는 (17.)로, 쌍대 공식은 (16.)으로 바뀐다. 최근 Müller는 Graßmann 전개의 이 특수한 경우들로부터\footnote{E. Müller, Beiträge zur Graßmannschen Ausdehnungslehre. 1. Mitteilung: Einige allgemeine Sätze. (Sitzungsber. d. Ak. d. Wiss. Wien; Bd. 118. Abt. IIa. Juli 1909), Satz 16 und 18.} 앞서 말한 항등식 (20.), (21.)을 유도하였다.

우리의 유도는 Graßmann--Müller의 유도와 달리 항등식의 완전성까지 동시에 증명한다. 여기서 다루는 문제에는 바로 이것이 가장 본질적이다. 또한 우리의 유도는 (20.), (21.)에서 출발하여 지금까지 주목되지 않은 듯한 가장 일반적인 분해 불가능 항등식 (36.)으로 나아간다.
